\section{When Groups said, "represent us!"}
Well, most often we do not see the group itself, but rather, its maniffestion. Like, say an element of the rotation group causes rotation of some object, due to which we are able to see its effects. To this extent, we define a group action.
\begin{definition}[Group Action]
    A group $(G,*)$ is said to act on some set $X$ if there exists a map $\phi: G\times X\rightarrow X$ such that:
    \begin{enumerate}
        \item $\phi(e,x) = x$, that is, action of identity element of group is to map the element of $X$ to itself.
        \item $\phi(g, \phi(h,x))= \phi(g*h, x)$, that is, action of a group element on another action is combination of the two actions.
    \end{enumerate}
\end{definition}
Group actions just cause permutations to the elements of $X$. Consider the example of $\re^3$, where the inner-product preserving group actions are just \textit{translations} and \textit{rotations}. \\[0.2cm]

Let $f$ be a function on $X$. Then the group action on the function is defined by:
$$(g\cdot f)(x) = f(g^{-1}x)$$
It is trivial that $e\cdot f (x) = f(x)$. Also, 
\begin{align*}
    g_1 \cdot (g_2 \cdot f)(x) &= g_2\cdot f(g_1^{-1}x)\\
    &=f(g_2^{-1}* g_1^{-1} x)\\
    &=f((g_1*g_2)^{-1}x)\\
    &=(g_1*g_2)f(x)
\end{align*}
Hence, using the definition as above, it is indeed a group action on the functions on $X$.
\begin{definition}[Representation]
    A \textit{representation} $(\pi,V)$ of a group $G$ is a group homomorphism defined as,
    $$\pi: G\rightarrow GL(V)$$
    where $GL(V)$ is the set of invertible linear maps $V\rightarrow V$ and $V$ is a vector space.
\end{definition}
Note that, if $V$ is finite-dimensional and a basis of $V$ is chosen, then the set of invertible linear maps can be seen in terms of matrices, which provide an isomorphism:
$$GL(V) \cong GL(n,\re)$$
where $GL(n,\re)$ is the group of $n\times n$ invertible, real matrices. In most cases, $G$ will be a Lie group and the representation $\pi$ will be a smooth map. \\[0.2cm]
There can be different representations of the group. The \textit{trivial} representation is such that $g\xmapsto{\pi}\mathrm{id}$, where $\mathrm{id}$ is the identity map. Given two representations $\pi_1$ and $\pi_2$ of dimensions $n_1$ and $n_2$ respectively, we can define the \textit{direct-sum} representation which is of dimension $n_1+n_2$ by the homomorphism:
$$(\pi_1\oplus \pi_2): G\rightarrow GL(V_{\pi_1}\oplus V_{\pi_2}) \quad \mathrm{s.t}\quad g\mapsto \mqty(\pi_1(g)& 0\\ 0&\pi_2(g))$$
This can be thought of as the representation matrix being a block diagonal matrix, with each sub-representation matrix occupying the diagonal blocks.
\begin{definition}[Irreducible Representation]
    An \textit{irreducible representation} (or irr rep.) $\pi$ is one which has no sub-representations, that is, $\nexists\ W\subset V$ such that the restriced function $\pi|_{W}$ is a representation.
\end{definition}
It turns out that if $\pi$ is an irreducible representation, it cannot be written as a direct-sum of two other representations. It makes sense, since direct-sum representation does not convey any new message and is not that interesting. Hence, looking at irreducible representations makes more sense.


\begin{lemma}[Schur's First Lemma]
        Let $\pi: G\rightarrow GL(V)$ and $\phi: G\rightarrow GL(U)$ be two irreducible representations of a finite group G. Then, if a matrix $A$ is such that:$$A\pi(g) = \phi(g)A \ \forall \ g\in G$$ then, either $A=0$ or $A$ is invertible (and hence two representations are equivalent).
\end{lemma}

\begin{lemma}[Schur's Second Lemma]
    Let $\pi$ be an irreducible, complex, representation of a group $G$ and $M$ is a linear map over the complex field, that commutes with $\pi(g)$ for all $g\in G$. Then, $M$ can atmost be a scalar multiplication with the identity, that is, $$M=\lambda \mathds{1},\ \lambda\in \ce$$
\end{lemma}
\textit{Proof.} Since $M$ is over the field of complex numbers, it must have atleast one eigenvalue which we call $\alpha$. Then $\det(M-\lambda\mathds{1}) = 0$ and hence $M-\lambda\mathds{1}$ is not invertible. \\[0.2cm]
Since $M\pi(g) = \pi(g)A$ for all $g$, we also have $(M-\lambda\mathds{1})\pi(g) = \pi(g)(M-\lambda\mathds{1})$ for all $g$. Then from Schur's first lemma, we have $M-\lambda\mathds{1} = 0$ since it is not invertible. Thus, $M=\lambda\mathds{1}$ for some $\lambda\in \ce$\\[0.2cm]
A corollolary of the above lemma says that if $G$ is a commutative group, then all of its irreducible representations will be one-dimensional.\\[0.2cm]
\textit{Proof.} Let $G$ be a commutative group and $\pi$ be its irreducible representation. Then, since group representation is a homomorphism, we have:
$$\pi(gh) = \pi(g)\pi(h) = \pi(hg) = \pi(h)\pi(g)\ \forall \ g,h\in G$$
The third step comes from the commutativity of the group elements. Now, fix $h$ and hence, for all $g$, we have the commutation relation holding true. This implies that $\pi(h) = \lambda_h \mathds{1}$ from second lemma. Since $h$ was arbitrary, all irreducible representations are one-dimensional $\pi(h) = \lambda_h \in \ce$.\\[0.2cm]
NOTE: The proof depends on existence of eigenvalues which will only be true for complex fields. For real vector spaces, this will in general not be true!

\subsection{Representations of $\boldsymbol{U(1)}$}
Consider rotations in 2D plane by some angle $\theta$. The most common entity we are reminded of is $e^{i\theta}$. It turns out that the group of rotations is a Lie group, denoted by $U(1)$, since those can be thought of as elements of a $1\times 1$ matrices.\\[0.2cm]
The elements of the group $U(1)$ are points on the unit-circle, labelled as $e^{i\theta}$ where $\theta$ can be any real number but $\theta$ and $\theta+2n\pi$ denote the same element, where $n\in \inte$. \\[0.2cm]
Since $U(1)$ is a commutative group, all its representations will be one-dimensional as seen before and is given by:
$$\pi: U(1)\rightarrow GL(1,\ce)$$
where $GL(1,\ce)$ is the group of all invertible complex numbers, that is, non-zero complex numbers. All the irreducible representations of $U(1)$ is unitary, that is, $\pi_k(\theta) = e^{ik\theta}, k\in \inte$
\subsubsection{Charge Operator and Symmetry}

